% Figure: the rectified flow model, in training and at sampling.
% Training mixes an image and noise at one random time on a straight line;
% sampling takes four Euler steps along the predicted direction.
% Facts come from the published model repository and model card.
% Requires in the preamble: \usepackage{tikz} + \usetikzlibrary{positioning,arrows.meta}
\begin{figure}[t]
  \centering
  \begin{tikzpicture}[
    x=1cm, y=1cm,
    font=\small,
    block/.style={rectangle, rounded corners=2pt, draw=black!70, line width=0.5pt,
      fill=blue!10, align=center, inner sep=4pt, minimum height=9mm},
    inputblock/.style={block, fill=green!12},
    stageblock/.style={block, fill=orange!15},
    noteblock/.style={rectangle, rounded corners=2pt, draw=black!45, dashed,
      fill=red!8, align=left, inner sep=3pt, font=\scriptsize},
    flow/.style={->, >=Latex, line width=0.7pt, draw=black!75},
    feed/.style={->, >=Latex, line width=0.7pt, draw=black!75, dashed},
    arrowlab/.style={font=\scriptsize, fill=white, inner sep=1pt},
  ]

  % ================= top half: how it was trained =================
  \node[anchor=west] (headT) at (0.15,10.45) {\textbf{How it was trained}};
  \node[font=\scriptsize, anchor=west, right=0.65cm of headT] (dirT)
    {one random time per example};

  % one straight line between the noise end and the image end
  \node[inputblock, anchor=west, text width=1.50cm, minimum height=1.15cm] (noiseT) at (0.15,9.20)
    {pure\\noise};
  \node[inputblock, anchor=east, text width=1.50cm, minimum height=1.15cm] (imageT) at (11.30,9.20)
    {clean\\image};
  \draw[flow] (noiseT.east) -- (imageT.west)
    coordinate[pos=0.16] (tp1) coordinate[pos=0.40] (tp2)
    coordinate[pos=0.64] (tp3) coordinate[pos=0.88] (tp4);
  \foreach \p in {tp1,tp3,tp4}{\fill[black!60] (\p) circle (1.4pt);}
  \fill[black] (tp2) circle (2.4pt);
  \node[font=\scriptsize, align=center, anchor=south] at ([yshift=3pt]tp2)
    {a mix at one\\random time $t$};
  \node[font=\scriptsize, anchor=south] at (8.15,9.27) {the learned direction};

  % the training loop: predict, compare, nudge, repeat
  \node[block, anchor=west, text width=3.00cm, minimum height=1.55cm] (lt1) at (0.15,7.35)
    {the model predicts the direction at that time};
  \node[block, text width=3.15cm, minimum height=1.55cm, right=0.40cm of lt1] (lt2)
    {compare the guess with the true direction along the line};
  \node[block, text width=1.75cm, minimum height=1.55cm, right=0.40cm of lt2] (lt3)
    {nudge the weights};
  \draw[feed] (tp2) -- ++(0,-0.70) -| (lt1.north);
  \draw[flow] (lt1) -- (lt2);
  \draw[flow] (lt2) -- (lt3);
  \draw[flow] (lt3.south) -- ++(0,-0.65) coordinate (retcT) -- (retcT -| lt1.south) -- (lt1.south);

  % ---------------- the one box that both halves use ----------------
  \path (lt1.center) -- (lt3.center) coordinate[midway] (loopmidT);
  \coordinate (retmidT) at (loopmidT |- retcT);
  \node[block, text width=3.40cm, minimum height=0.90cm] (net) at ([yshift=-0.90cm]retmidT)
    {\textbf{the same network}};
  \draw[feed] (retmidT) -- (net.north);

  % ================= bottom half: what we actually run =================
  \node[inputblock, anchor=west, text width=1.35cm, minimum height=1.15cm] (noiseB) at (0.15,3.65)
    {pure\\noise};
  \node[stageblock, text width=2.20cm, minimum height=1.35cm, right=0.45cm of noiseB] (st1)
    {\textbf{Euler step 1}\\[1pt]\scriptsize step along\\the direction};
  \node[stageblock, text width=2.20cm, minimum height=1.35cm, right=0.45cm of st1] (st2)
    {\textbf{Euler step 2}\\[1pt]\scriptsize step along\\the direction};
  \node[stageblock, text width=2.20cm, minimum height=1.35cm, right=0.45cm of st2] (st3)
    {\textbf{Euler step 3}\\[1pt]\scriptsize step along\\the direction};
  \node[stageblock, text width=2.20cm, minimum height=1.35cm, right=0.45cm of st3] (st4)
    {\textbf{Euler step 4}\\[1pt]\scriptsize step along\\the direction};
  \node[inputblock, text width=1.40cm, minimum height=1.15cm, right=0.45cm of st4] (latentB)
    {clean\\latent};
  \draw[flow] (noiseB) -- (st1);
  \draw[flow] (st1) -- (st2);
  \draw[flow] (st2) -- (st3);
  \draw[flow] (st3) -- (st4);
  \draw[flow] (st4) -- (latentB);

  % the same network drives every Euler step
  \coordinate (busY) at (0,4.46);
  \draw[flow] (net.south) -- (net.south |- busY);
  \draw[draw=black!75, line width=0.7pt] (st1.north |- busY) -- (st4.north |- busY);
  \foreach \s in {st1,st2,st3,st4}{\draw[flow] (\s.north |- busY) -- (\s.north);}

  \node[anchor=west] (headB) at (0.15,2.35) {\textbf{What we actually run}};
  \node[font=\scriptsize, anchor=west, right=0.65cm of headB] (dirB) {four Euler steps};

  % ================= notes =================
  \node[noteblock, text width=3.55cm, anchor=north west] at (11.50,10.35)
    {Training needs a very large image set and many GPU hours. We trained nothing: the
     base checkpoint and the community-trained adapter are both released weights.};
  \node[noteblock, text width=6.95cm, anchor=west] at (7.70,5.40)
    {The model predicts the direction, not a noise. Every Euler step moves the sample
     along that direction.};
  \node[noteblock, text width=7.20cm, anchor=north west] at (0.15,1.85)
    {The model is step distilled, so four Euler steps are close to what a slow model would
     reach in dozens of steps.};
  \node[noteblock, text width=7.20cm, anchor=north west] at (7.65,1.85)
    {We never run the training path to make an image. It only exists so the model can learn.};

  \end{tikzpicture}
  \caption{The rectified flow model in training and at sampling. Training draws one random
    time per example, mixes the image and noise at that time, and teaches the model the
    direction along the straight line between the two ends. Sampling starts from noise and
    takes four Euler steps along the predicted direction. Both jobs use one and the same
    network.}
  \label{fig:rectified_flow_process}
\end{figure}
